\documentclass[11pt]{article}
\usepackage[a4paper,margin=27mm]{geometry}
\usepackage{amsmath,amssymb,amsthm,mathtools}
\usepackage{graphicx}
\usepackage{microtype}
\usepackage[hidelinks]{hyperref}

\newtheorem{theorem}{Theorem}
\newtheorem{proposition}{Proposition}
\newtheorem{example}{Example}
\newcommand{\N}{\mathbb N}
\newcommand{\C}{\mathbb C}
\newcommand{\D}{\mathcal D}

\title{A no-go theorem for the literal unbounded-operator extension\\
\large Fractional difference equations with finite delay}
\author{Candidate partial-result packet}
\date{August 11, 2026}

\begin{document}
\maketitle

\noindent\textbf{Status.}
Candidate sharp obstruction, likely valid pending expert review.

\section{Source direction}

Zhang and Bu study, for $2<\alpha<3$, the explicit fractional difference
equation
\begin{equation}\label{eq:model}
 \Delta^\alpha u(n)=Au(n)+\gamma u(n-\lambda)+f(n),
 \qquad n\in\N_0,
\end{equation}
with
\begin{equation}\label{eq:initial}
 u(-\lambda)=\cdots=u(0)=u(1)=u(2)=0.
\end{equation}
Their coefficient $A$ is bounded. The introduction says that systems
including unbounded operators remain an open problem
\cite{ZhangBu}. The source statement is a broad research direction rather
than a quantified conjecture.

\begin{figure}[ht]
  \centering
  \includegraphics[width=\textwidth]{figures/open_problem_crop.png}
  \caption{The unbounded-operator direction in the source introduction,
  PDF page 3.}
\end{figure}

This packet tests the closest literal extension: let
$A:\D(A)\subseteq X\to X$ be closed, keep \eqref{eq:model}--\eqref{eq:initial},
and ask for a classical solution for every finitely supported $X$-valued
forcing. A classical solution must satisfy $u(n)\in\D(A)$ whenever
$Au(n)$ occurs. This is the solution class needed by the usual maximal
regularity requirement $Au\in\ell^p(\N_0;X)$.

\section{Proof intuition}

The scheme is explicit: the third forward difference at time zero determines
the value three time levels ahead. Because the first three values of $u$ are
zero, an impulse $f(0)=x$ passes through without smoothing and gives
$u(3)=x$. Three steps later, the same equation asks to apply $A$ to $u(3)$.
Thus every admissible impulse must already lie in $\D(A)$. If all
$X$-valued impulses are allowed, then $\D(A)=X$, and closedness makes $A$
bounded.

This same obstruction is encoded in the source's resolvent sequence:
$S_\alpha(0)=I$, so the first recurrence step contains $AI=A$. Retaining a
bounded-operator-valued resolvent sequence also forces boundedness. Finally,
a diagonal example shows that the frequency-domain resolvent bounds alone
cannot restore classical solvability.

\section{The domain obstruction}

The source defines
\[
 \Delta^\alpha=\Delta^3\Delta^{-(3-\alpha)},\qquad
 \Delta^{-(3-\alpha)}u=k^{3-\alpha}*u,
\]
where $k^\beta(0)=1$ and $\Delta v(n)=v(n+1)-v(n)$.

\begin{theorem}[All-data rigidity]\label{thm:rigidity}
Let $X$ be a Banach space, $2<\alpha<3$, $\lambda\in\N$, and
$\gamma\in\C$. Let $A:\D(A)\subseteq X\to X$ be a linear operator.
Suppose that for every finitely supported $f:\N_0\to X$ there is a classical
solution of \eqref{eq:model}--\eqref{eq:initial} on all of $\N_0$. Then
$\D(A)=X$. In particular, if $A$ is closed, then $A\in\mathcal B(X)$.
\end{theorem}

\begin{proof}
Fix an arbitrary $x\in X$ and take the impulse forcing
\[
 f(0)=x,\qquad f(n)=0\quad(n\ge1).
\]
Set $v=\Delta^{-(3-\alpha)}u=k^{3-\alpha}*u$. From
\eqref{eq:initial},
\[
 v(0)=v(1)=v(2)=0,\qquad v(3)=k^{3-\alpha}(0)u(3)=u(3).
\]
Consequently,
\[
 \Delta^\alpha u(0)=\Delta^3v(0)
 =v(3)-3v(2)+3v(1)-v(0)=u(3).
\]
Evaluating \eqref{eq:model} at $n=0$ and using the zero initial history gives
\[
 u(3)=Au(0)+\gamma u(-\lambda)+f(0)=x.
\]
Since the solution is classical on all of $\N_0$, equation
\eqref{eq:model} at $n=3$ requires $u(3)\in\D(A)$. Hence $x\in\D(A)$.
The vector $x\in X$ was arbitrary, so $\D(A)=X$. A closed everywhere
defined linear operator on a Banach space is bounded by the closed graph
theorem.
\end{proof}

\begin{proposition}[The source resolvent sequence is rigid]\label{prop:S}
Suppose one keeps Definition 3.1 of \cite{ZhangBu}: all
$S_\alpha(n)$ are bounded operators,
$S_\alpha(0)=S_\alpha(1)=S_\alpha(2)=I$,
$S_\alpha(-j)=0$ for $1\le j\le\lambda$, and the displayed recurrence is an
identity in $\mathcal B(X)$. Then its coefficient $A$ is bounded. Thus this
definition has no genuinely unbounded instance.
\end{proposition}

\begin{proof}
At $n=0$,
\[
 (k^{\alpha-2}*S_\alpha)(0)=S_\alpha(0)=I,\qquad
 (k^{\alpha-2}*S_\alpha^\lambda)(0)=S_\alpha(-\lambda)=0.
\]
The recurrence therefore expresses $A$ as $S_\alpha(3)$ plus a scalar
multiple of $I$. Both are bounded.
\end{proof}

\section{The frequency bounds do not rescue the extension}

Write
\[
 g(z)=z^{3-\alpha}(z-1)^\alpha,\qquad
 q(z)=g(z)-\gamma z^{-\lambda}\qquad(|z|=1).
\]
The source's characterization uses the two families
\begin{equation}\label{eq:families}
 \{g(z)(q(z)-A)^{-1}\},\qquad
 \{z^{-\lambda}(q(z)-A)^{-1}\},
\end{equation}
away from the exceptional points of the circle.

\begin{example}[Frequency conditions without classical solvability]
\label{ex:diag}
Let $H=\ell^2(\N;\C)$, put $M=2^\alpha+|\gamma|$, and define
\[
 Ae_j=-(M+2+j)e_j,\qquad
 \D(A)=
 \left\{x=(x_j):\sum_{j\ge1}(M+2+j)^2|x_j|^2<\infty\right\}.
\]
Then $A$ is densely defined, self-adjoint, closed, and unbounded. Both
families in \eqref{eq:families} are $R$-bounded, but
\eqref{eq:model}--\eqref{eq:initial} is not classically solvable for every
finitely supported $H$-valued forcing.
\end{example}

\begin{proof}
For $|z|=1$,
\[
 |q(z)|\le |z-1|^\alpha+|\gamma|\le 2^\alpha+|\gamma|=M.
\]
Since $\sigma(A)=\{-(M+2+j):j\ge1\}$,
\[
 \operatorname{dist}(q(z),\sigma(A))\ge3,\qquad
 \|(q(z)-A)^{-1}\|\le\frac13.
\]
Also $|g(z)|\le2^\alpha$. Hence both families in
\eqref{eq:families} are uniformly norm bounded. On a Hilbert space,
uniform norm boundedness is equivalent to $R$-boundedness.

Now take $x=(1/j)_{j\ge1}$. Then $x\in H$ but $x\notin\D(A)$. The impulse
$f(0)=x$, $f(n)=0$ for $n\ge1$, belongs to every
$\ell^p(\N_0;H)$. The proof of Theorem~\ref{thm:rigidity} forces
$u(3)=x$, which is outside $\D(A)$, so there is no classical solution on all
of $\N_0$.
\end{proof}

\section{Consequences and limitations}

There are two exhaustive outcomes for a naive replacement of
$A\in\mathcal B(X)$ by a closed operator while retaining the source
hypotheses:

\begin{itemize}
\item If the bounded source-style resolvent sequence is retained, then
Proposition~\ref{prop:S} forces $A$ to be bounded, so the extension has no new
operators.
\item If that hypothesis is dropped and only the frequency families are
retained, Example~\ref{ex:diag} shows that they do not imply all-data
classical solvability.
\end{itemize}

Allowing a mild equation in an extrapolation space does not restore the
standard conclusion $Au\in\ell^p(\N_0;X)$ for arbitrary $X$-valued data:
the impulse still gives $u(3)=x$, and $Au(3)\in X$ requires $x\in\D(A)$.

This is a complete no-go result for the literal extension, but only a partial
answer to the source's broad research direction. A positive unbounded theory
must materially reformulate the problem---for example by imposing compatible
data in a domain scale, weakening the target regularity, or using an implicit
scheme where a resolvent of $A$ smooths the newest time level.

\section{Verification and novelty status}

The time-zero identity was independently derived from the source definition
and cross-checked against the same computation in its proof of Theorem 3.8.
The diagonal example was checked by its exact spectrum and the elementary
symbol bound. No numerical evidence is used.

On August 11, 2026, the local result indexes and bounded exact-title,
quoted-phrase, and core-term web searches found no later resolution of this
obstruction. Novelty confidence is moderate pending a specialist citation
search.

\begin{thebibliography}{9}

\bibitem{ZhangBu}
J. Zhang and S. Bu,
\emph{Maximal regularity for fractional difference equations with finite
delay on UMD spaces},
arXiv:2406.15417 (2024).

\bibitem{Blunck}
S. Blunck,
\emph{Maximal regularity of discrete and continuous time evolution
equations},
Studia Math. \textbf{146} (2001), 157--176.

\end{thebibliography}

\end{document}
